\subsection{Pullback Bundles}

    Given a fiber bundle $F \hookrightarrow E \oset[0.7pt]{\scalebox{0.8}{$p$}}{\longrightarrow} B$ and a map $g \colon B^\prime \to B$ we can construct the pullback bundle $g^*E$ using Theorem \ref{thm:fiber_bundle_construction2} as the fiber bundle over $B^\prime$, where the fiber over $b^\prime \in B^\prime$ is given by the fiber of $E$ over $g(b^\prime)$, i.e.
    \[
    g^*E = \coprod_{b^\prime \in B^\prime} p^{-1}(g(b^\prime)).
    \]
    Given a local trivialisation $\phi \colon U \times F \to E$ we define a local trivialisation of $g^*E$ by
    \[
    g^{-1}(U) \times F \to g^*E, \; (b^\prime,f) \mapsto (b^\prime,\phi(g(b^\prime),f)).
    \]
    The same construction works in the smooth category using Theorem \ref{thm:fiber_bundle_construction2_smooth} as long as either $F$ has no boundary or $B$ and $B^\prime$ have no boundary. By projecting onto the second factor we get a bundle homomorphism $G \colon g^*E \to E$, making the following diagram commute:
    \[
        \commsquare{$g^*E$}{$E$}{$B^\prime$}{$B$}{$G$}{}{$p$}{$g$}
    \]

    \begin{lemma}
        \label{lem:pullback_of_bundle_under_embedding}
        In the above situation the following holds:
        \begin{enumerate}[(i)]
            \item If $g$ is an immersion, so is $G$
            \item If $g$ is an embedding, so is $G$
        \end{enumerate}
    \end{lemma}
   \begin{proof}
    In local trivialisations we can write $G$ as
    \[
    g^{-1}(U) \times F \to U \times F, \; (b,f) \mapsto (g(b^\prime),f),
    \]
    such that (i) follows. Now let $g$ be an embedding. We have $\mathrm{im} G = p^{-1}(g(B^\prime))$, so we can define 
    \[
    H \colon \mathrm{im} G \to f^*E, \; x \mapsto (g^{-1}(p(x)),x).
    \]
    Clearly $H$ is continuous and $G \circ H = \mathrm{id}_{\mathrm{im} G}$ as well as $H \circ G = \mathrm{id}_{g^*E}$, so $G$ is in fact an embedding.
   \end{proof}

   Given a fiber bundle $F \hookrightarrow E \oset[0.7pt]{\scalebox{0.8}{$p$}}{\longrightarrow} B$ and a submanifold $S \subseteq B$ we can then define the restriction of $E$ to $S$ as the pullback of $E$ under the inclusion $\iota \colon S \hookrightarrow B$:
   \[
   E|_S := \iota^*E
   \]
   The above Lemma shows that this is indeed a submanfold of $E$. In fact, in the situation of the above Lemma we have $$\mathrm{im} G = E|_{g(B^\prime)}.$$

   \vspace{1em}

   One usecase for pullback bundles is to define vector fields along curves $\gamma \colon I \to M$ for a smooth manifold $M$ as sections in the pullback bundle $\gamma^\ast TM$. The follwoing Lemma says, that we can equivalently define them as maps 
   \[
   V \colon I \to TM
   \]
   with $V(t) \in T_{\gamma(t)} M$ (as is done in \cite{lee2018riemannian}).

   \begin{lemma}
    Let $F \hookrightarrow E \to M$ be a smooth fiber bundle and $g \colon N \to M$ a smooth map. Then a section $N \xrightarrow{\sigma} g^\ast E$ is smooth if and only if the composition
    \[
    N \xrightarrow{\sigma} g^\ast E \xrightarrow{G} E
    \]
    is smooth, where $G$ is the induced bundle homomorphism $g^\ast E \to E$ discussed above.
   \end{lemma}

   \begin{proof}
    This follows directly from the definition of the smooth structure of the pullback bundle given above.
   \end{proof}